% !TeX root = ../main.tex
\begin{abstract}
Pooled CRISPR screens increasingly use longitudinal, donor-adjusted, and
factorial designs, but beta-binomial screen methods have largely remained
limited to pairwise comparisons.  BARCS extends the
library-total-conditional beta-binomial model to guide-level regression with
an arbitrary design matrix, so that time, covariate, and interaction effects
can be estimated in a single fit.  In four replicate-complete Cas13 screens,
adding the intermediate time point gave a modest improvement in
essential-gene recovery.  Applying the same non-targeting-control scaling
rule to BARCS, MaGeCK-MLE, edgeR-QL, DESeq2, and limma-voom gave similar
calibration across all five methods, though the four alternatives ranked
essential genes more strongly than BARCS.  In an ordered-bin IL2RA screen,
donor-adjusted BARCS recovered more validated regulators with fewer total
calls than the matched four-bin MAGeCK-MLE fit, and cross-fitted controls
exposed excess guide-level significance.  Simulations showed gains from
dispersion moderation and control-based denominators, but seed-specific
results exposed sensitivity to the denominator choice, and a null grid
traced the large gene-level error to correlated-guide aggregation rather
than to dispersion alone.  Aggregation-matched control scaling reduced this
error but did not remove it.  An external audit of a published beta-binomial
screen showed that the reported CB\textsuperscript{2}
null-discovery count disappeared once full-library totals were restored.
That corrected one denominator-dependent result but did not resolve the
broader calibration concern; nominal-level calibration remains unresolved.
BARCS is a multivariable extension of the beta-binomial CRISPR model, not a
claim of uniform superiority over negative-binomial methods.  Confirmatory
use requires independent biological replication, likelihood-compatible
counts, and held-out or cross-fitted controls.
\end{abstract}
